% Pista geo-26j2-q4b · Geometria
% Procedència: PAU Catalunya 2026 · Convocatòria de juny · Sèrie 5 · Exercici 4 — Opció B
\documentclass[11pt,a4paper]{article}
\usepackage{pista}
\pistainfo{Catalunya 2026}{Convocatòria de juny}{Sèrie 5}

\begin{document}

\section*{\textcolor{blauPista}{Exercici 4 --- Opció B}}

\noindent En una indústria es fan servir braços robòtics amb sensors per a mesurar amb precisió determinades peces metàl·liques i comprovar si tenen irregularitats. Les peces es col·loquen sobre la superfície determinada pel pla $\pi: (x, y, z) = (6, 1, 1) + \lambda(1, 1, 0) + \mu(0, 0, 1)$ i, per damunt seu, un braç robòtic desplaça un sensor des del punt $A = (1, 2, 3)$ fins al punt $B = (5, 3, 7)$ en línia recta (les unitats dels tres eixos de coordenades estan expressades en centímetres).

\subsection*{4.1. \normalfont Quina distància ha recorregut el sensor? \hspace*{\fill}\textnormal{[0,5~punts]}}

\pista{La distància que ha recorregut el sensor és el mòdul del vector $\vec{AB}$. Recorda que el mòdul d'un vector $\vec{v} = (v_{1}, v_{2}, v_{3})$ es calcula així:
\[
|\vec{v}| = \sqrt{v_{1}^{2} + v_{2}^{2} + v_{3}^{2}}
\]
Les unitats són centímetres.}

\subsection*{4.2. \normalfont En acabar el moviment, a quina distància de la superfície es troba el sensor? \hspace*{\fill}\textnormal{[1~punt]}}

\pista{El pla $\pi$ et ve donat amb un punt i dos vectors directors. En primer lloc, expressa el pla amb l'equació general: \mbox{$ax + by + cz + d = 0$}. Per aconseguir-ho, iguala a zero aquest determinant, on a les files segona i tercera hi van els dos vectors directors:
\[
\begin{vmatrix} x - 6 & y - 1 & z - 1 \\ \ldots & \ldots & \ldots \\ \ldots & \ldots & \ldots \end{vmatrix} = 0
\]
A continuació, aplica la fórmula de la distància d'un punt $(x_{0}, y_{0}, z_{0})$ a un pla al punt $B$, que és on acaba el sensor:
\[
d(B, \pi) = \frac{|a x_{0} + b y_{0} + c z_{0} + d|}{\sqrt{a^{2} + b^{2} + c^{2}}}
\]}

\subsection*{4.3. \normalfont Si el sensor seguís movent-se en la mateixa línia recta i en el mateix sentit, en quin punt xocaria amb la superfície? \hspace*{\fill}\textnormal{[1~punt]}}

\pista{El punt on xocaria és la intersecció de la recta $AB$ amb el pla $\pi$. En primer lloc, escriu la recta en forma paramètrica, amb el punt $A$ i el vector director $\vec{AB}$. Fixa't que quan $\lambda = 0$ obtens el punt $A$, i que quan $\lambda = 1$ obtens el punt $B$.

\medskip
A continuació, substitueix les expressions de $x$, $y$ i $z$ a l'equació del pla de l'apartat 4.2: obtindràs una equació amb la variable $\lambda$.

\medskip
Un cop l'hagis resolta, substitueix el valor de $\lambda$ a l'equació de la recta. Com que el sensor continua endavant, el valor de $\lambda$ ha de ser més gran que 1, i això et permetrà escollir la solució que és vàlida en aquest context.}

\end{document}
